% ============================================================
%  Chapter 16 — Cooperative Composition and Gift Exchange
% ============================================================
\chapter{Cooperative Composition and Gift Exchange}
\label{ch:cooperative}

\begin{quote}
\textit{``The gift not yet repaid debases the man who accepted it.''}\\
--- Marcel Mauss, \textit{The Gift} (1925)
\end{quote}

\vspace{1em}

\section{Introduction}

Marcel Mauss's celebrated \textit{Essai sur le don} (1925) demonstrated that
gift exchange is never merely economic: to give is to create obligation, to
receive is to assume debt, and to reciprocate is to weave a social bond.
In societies from the Trobriand Islands to the Pacific Northwest, the
circulation of gifts constitutes a total social fact---simultaneously
economic, juridical, moral, and aesthetic.\footnote{Mauss distinguished three
obligations of the gift: to give, to receive, and to reciprocate. Each is
indispensable; failure at any stage ruptures the social fabric.}

In Ideatic Diffusion Theory, cooperation is \emph{composition}: when two agents
contribute their ideas to a joint project, the result is $\compose(a, b)$.
The order matters---gifts are not symmetric!---and the void $\void$
contributes nothing.  This chapter formalizes cooperative game theory within
ideatic space, establishing twenty theorems that characterize joint
composition, multi-party projects, depth bounds on cooperative output, and
resonance preservation under collaboration.

The analogy between $\compose$ and gift-giving is precise.  Agent $a$
``gives first,'' setting the frame of the exchange; agent $b$ reciprocates
within that frame.  The result $\compose(a,b)$ is a compound idea whose
depth is bounded by the sum of the contributors' depths---one cannot create
profundity from banality---and whose resonance properties are inherited from
the inputs.  When we extend to multi-party projects via iterated
composition, we obtain a formal model of Malinowski's \textit{Kula ring}, of
potlatch sequences, and of any process in which meaning accrues through
sequential cooperative acts.\footnote{Bronislaw Malinowski's
\textit{Argonauts of the Western Pacific} (1922) documented how shell
valuables circulate through a ring of trading partners, each exchange adding
a layer of social meaning.}

We proceed as follows.  Section~\ref{sec:ch16-defs} introduces the core
definitions---\texttt{jointCompose}, \texttt{symmetricJoint}, and
\texttt{JointProject}---together with their Lean formalizations.
Section~\ref{sec:ch16-void} establishes the role of void as the
non-contributor.  Section~\ref{sec:ch16-assoc} proves associativity of
multi-party cooperation.  Section~\ref{sec:ch16-resonance} explores
resonance preservation.  Section~\ref{sec:ch16-depth} derives depth
bounds on cooperative output.  Section~\ref{sec:ch16-multi} extends the
analysis to multi-party and iterated cooperation.  Section~\ref{sec:ch16-concl}
concludes.

\subsection*{Mauss's Three Obligations and the Structure of Reciprocity}

Mauss's analysis of gift exchange identified three fundamental
obligations: the obligation \emph{to give}, the obligation \emph{to
receive}, and the obligation \emph{to reciprocate}.  Each obligation is
indispensable; failure at any stage ruptures the social bond.  In our
formal framework, these three obligations map onto structural properties
of composition.  The obligation to give corresponds to the act of
contributing $a$ to a joint composition $\compose(a, b)$---the initiator
must provide an input.  The obligation to receive corresponds to the
requirement that the second argument $b$ accept the compositional frame
established by $a$---compositionally, $b$ operates within the context
set by $a$, which it cannot refuse without leaving the exchange entirely.
The obligation to reciprocate corresponds to the non-triviality of $b$:
reciprocation with $\void$ (Theorem~\ref{thm:void_joint_right}) leaves
the initiator's gift unchanged, which is structurally equivalent to
non-reciprocation.  A genuine return gift requires $b \neq \void$, so
that $\compose(a, b) \neq a$.

This three-fold structure reveals why gift exchange creates lasting
social bonds in a way that market transactions do not.  A market
transaction terminates at the moment of exchange: money for goods, and
the obligation is discharged.  A gift, by contrast, creates an
asymmetric temporal structure: the obligation to reciprocate persists
until a counter-gift is offered, and the counter-gift itself creates a
new obligation, generating an indefinitely extended chain of mutual
obligation.  In compositional terms, the gift exchange generates an
iterated sequence $\compose(a_1, \compose(b_1, \compose(a_2,
\compose(b_2, \ldots))))$ that is structurally open-ended---there is no
natural termination point, because each reciprocation is itself a gift
that demands reciprocation.\footnote{This open-endedness distinguishes
gift exchange from contract: a contract specifies termination conditions,
while a gift relationship is, in principle, interminable.  See Marshall
Sahlins, \textit{Stone Age Economics} (1972), Chapter 5.}

\subsection*{Sahlins's Typology of Reciprocity}

Marshall Sahlins, in \textit{Stone Age Economics} (1972), proposed a
typology of reciprocity that maps naturally onto our compositional
framework.  \emph{Generalized reciprocity} is the altruistic extreme:
$a$ gives without expectation of return, corresponding to
$\compose(a, \void)$---the gift is offered and the return is void.
\emph{Balanced reciprocity} involves equivalent counter-gifts:
$\compose(a, b)$ where $\depth(b) \approx \depth(a)$, so that the
exchange is approximately symmetric in depth.  \emph{Negative
reciprocity} is the exploitative extreme: one party attempts to extract
maximum value while giving minimum in return, corresponding to
$\compose(a, b)$ where $\depth(a) \gg \depth(b)$ or $\depth(b) \gg
\depth(a)$, and the asymmetry is intentional rather than freely chosen.

Sahlins arranged these types on a continuum correlated with social
distance: generalized reciprocity prevails among close kin (parents
give to children without expecting repayment), balanced reciprocity
governs relations among trading partners and allies, and negative
reciprocity characterizes dealings with strangers and enemies.  Our
framework captures this continuum through the depth ratio
$\depth(a)/\depth(b)$: values near 1 indicate balanced exchange,
values much greater or less than 1 indicate asymmetric exchange, and
the direction of asymmetry (which party contributes more depth)
determines whether the exchange is generous or exploitative.

\subsection*{Gregory, Weiner, and the Inalienability Problem}

Chris Gregory's \textit{Gifts and Commodities} (1982) drew a sharp
distinction between gift economies and commodity economies.  In a
commodity economy, the exchange transfers \emph{alienable} objects
between \emph{independent} parties.  In a gift economy, the exchange
transfers \emph{inalienable} possessions between \emph{interdependent}
parties.  The key difference is that the gift retains a connection to
its giver: the \textit{hau} (spirit) of the gift, in Maori terms,
compels reciprocation.

Annette Weiner, in \textit{Inalienable Possessions} (1992), extended
this analysis by identifying a paradoxical practice she called
``keeping-while-giving.''  Certain treasured possessions---royal
regalia, sacred heirlooms, clan emblems---circulate as gifts but are
never fully alienated from their original owners.  The giver retains
a claim on the object even after giving it away.  In our compositional
framework, this corresponds to an idea $a$ that participates in
composition $\compose(a, b)$ without being ``consumed'': $a$ remains
available for future compositions even after its contribution to the
current joint project.  This is structurally guaranteed by our model,
since composition does not destroy its inputs---$a$ and $b$ continue
to exist as elements of $\IS$ after $\compose(a, b)$ is computed.

The inalienability of the gift maps onto the persistence of ideas in
ideatic space.  When a scholar shares an idea through publication, the
idea enters the public repertoire without leaving the scholar's
repertoire---this is ``keeping-while-giving'' in the academic context.
The formal framework captures this naturally: expansion adds an idea
to the recipient's repertoire without removing it from the giver's.

\subsection*{Modern Gift Exchange: From Blood to Code}

The logic of gift exchange extends far beyond pre-modern societies.
Richard Titmuss's \textit{The Gift Relationship} (1970) analyzed blood
donation as a form of generalized reciprocity: donors give blood to
strangers, expecting no direct return, motivated by solidarity with an
abstract community.  In our terms, each blood donation is a contribution
$a_{\text{blood}}$ to a joint project
$\texttt{output}(\langle [a_1, a_2, \ldots, a_n] \rangle)$ whose
output is the functioning of the healthcare system.  By
Theorem~\ref{thm:empty_project_void}, the healthcare system without
donors produces $\void$---nothing.  Each donor's contribution is
indispensable in the aggregate.

Open-source software represents another form of cooperative
composition.  When a programmer contributes code to a public
repository, they compose their work with that of other contributors:
$\compose(\text{existing\_codebase}, \text{new\_contribution})$.  The
Linux kernel, with its thousands of contributors over three decades,
is a monumental joint project whose output is the iterated composition
of individual contributions.  By
Theorem~\ref{thm:iterated_self_cooperation_depth}, the depth of this
project is bounded by the sum of individual depths---a bound that, for
Linux, runs into the millions of person-years of intellectual investment.

Jacques Derrida's \textit{Given Time} (1991) posed a radical
philosophical challenge to the very concept of the gift.  Derrida
argued that a true gift---one given without any expectation of return,
without even the consciousness that one is giving---is structurally
impossible within an economy of exchange.  The moment the gift is
recognized as a gift, it enters the economy of obligation and ceases
to be gratuitous.  In our framework, Derrida's argument translates to
the observation that $\compose(a, b)$ is always structurally
\emph{responsive}: the second argument $b$ is shaped by $a$, and the
result $\compose(a, b)$ bears the trace of both contributions.  A
``pure gift'' would require $\compose(a, b) = b$---that the giver's
contribution leave no trace---which is precisely the void condition
$a = \void$.  Derrida's paradox is thus that the only true gift is
no gift at all.\footnote{Lewis Hyde's \textit{The Gift} (1983) offers
a less paradoxical account, arguing that creative gifts---artworks,
stories, scientific discoveries---circulate through communities in ways
that create bonds of gratitude and obligation.  Hyde's ``gift economy''
is closer to Mauss's balanced reciprocity than to Derrida's impossible
ideal.}

% ============================================================
\section{Core Definitions}
\label{sec:ch16-defs}

We begin with the fundamental structures of cooperative composition.

\begin{definition}[Joint Composition]
\label{def:jointCompose}
Let $a, b \in \IS$.  The \emph{joint composition} of $a$ and $b$ is
\[
  \texttt{jointCompose}(a, b) \;:=\; \compose(a, b).
\]
In Mauss's terms, agent $a$ gives first and agent $b$ reciprocates.
The result is a shared compound idea.
\end{definition}

\begin{lstlisting}[caption={Joint composition in Lean}]
def jointCompose (a b : I) : I := IdeaticSpace.compose a b
\end{lstlisting}

Joint composition is simply composition itself, but the naming emphasizes the
cooperative framing: two agents are contributing sequentially to a shared
outcome.  The asymmetry is essential.  In Mauss's analysis, the initiator of
an exchange determines its moral character; the recipient responds within a
framework already established.

\begin{definition}[Symmetric Joint]
\label{def:symmetricJoint}
Two ideas $a, b$ are \emph{symmetrically joint} when the order of
composition does not matter:
\[
  \texttt{symmetricJoint}(a, b) \;:\Longleftrightarrow\;
  \compose(a, b) = \compose(b, a).
\]
\end{definition}

\begin{lstlisting}[caption={Symmetric joint in Lean}]
def symmetricJoint (a b : I) : Prop :=
  IdeaticSpace.compose a b = IdeaticSpace.compose b a
\end{lstlisting}

Symmetric joints are rare.  In most real exchanges---ritual, economic,
discursive---who gives first shapes the meaning of the entire exchange.
The symmetric case corresponds to perfectly reciprocal partnerships where
neither party holds agenda-setting power.

\begin{definition}[Joint Project]
\label{def:JointProject}
A \emph{joint project} consists of a list of participants, each contributing
an idea.  The project's \emph{output} is the left-fold of all contributions
via composition, starting from $\void$:
\[
  \texttt{output}(p) \;:=\; \text{foldl}(\compose, \void, p.\text{participants}).
\]
The \emph{total depth} of a joint project is $\sum_{a \in p} \depth(a)$.
\end{definition}

\begin{lstlisting}[caption={Joint project structure in Lean}]
structure JointProject (I : Type*) [IdeaticSpace I] where
  participants : List I

def JointProject.output (p : JointProject I) : I :=
  p.participants.foldl IdeaticSpace.compose IdeaticSpace.void

def JointProject.totalDepth (p : JointProject I) : N :=
  depthSum p.participants
\end{lstlisting}

The left-fold semantics mean that the first contributor sets the initial
frame (after void), and each subsequent contributor composes onto the
accumulating result.  This models the sequential accretion of meaning in
collaborative enterprises---from barn-raisings to co-authored papers.

\begin{example}[The Kula Ring as Iterated Cooperative Composition]
\label{ex:kula-ring}
The Kula Ring, documented by Bronislaw Malinowski in \textit{Argonauts
of the Western Pacific} (1922) and further analyzed by Nancy Munn in
\textit{The Fame of Ganda} (1986), is among the best-studied systems of
ceremonial gift exchange.  In the Kula, shell valuables (\textit{vaga},
arm-shells, and \textit{mwali}, necklaces) circulate in opposite
directions through a ring of trading partners spanning the islands of
eastern Papua New Guinea.

In our formal framework, the Kula Ring is an iterated joint project.
Let $a_1, a_2, \ldots, a_n$ represent the traders arranged in a ring,
and let $g_i$ represent the gift offered by trader $a_i$.  The
composition $\compose(g_1, \compose(g_2, \ldots, \compose(g_{n-1},
g_n)\ldots))$ represents the accumulated social meaning of the exchange
cycle.  By Theorem~\ref{thm:gift_triangle_depth} (generalized to $n$
parties), the depth of this composed meaning is bounded by
$\sum_{i=1}^n \depth(g_i)$.

What makes the Kula remarkable, from our formal perspective, is that
the valuables themselves gain depth through circulation.  Munn showed
that a shell's \emph{fame}---its social value---increases with each
exchange, because the history of its previous owners becomes part of
its meaning.  In IDT terms, the gift $g_i$ at step $i$ is not the
same as the ``raw'' shell but $\compose(g_{i-1}, g_i)$: the shell
\emph{composes with} each new exchange context.  The bounded depth of
each step (by subadditivity) means that the shell's accumulated meaning
grows at most linearly with the number of exchanges---but it does grow,
and this growth is precisely what gives old, much-traveled shells their
extraordinary prestige.
\end{example}

\begin{example}[Blood Donation as Generalized Reciprocity]
\label{ex:titmuss}
Richard Titmuss's \textit{The Gift Relationship} (1970) compared
voluntary (British) and commercial (American) blood donation systems,
arguing that the voluntary system produced higher-quality blood,
greater supply reliability, and stronger social cohesion.  In IDT
terms, the voluntary system models \emph{generalized reciprocity}
(Sahlins's typology): each donor $a_i$ contributes to a joint project
$\texttt{output}(\langle [a_1, \ldots, a_n] \rangle)$ without knowing
or caring about the identity of the recipient.

The key formal property is that generalized reciprocity requires
$b = \void$ in the donor's local exchange: the donor gives
($\compose(a_{\text{donor}}, \void) = a_{\text{donor}}$ by
Theorem~\ref{thm:void_joint_right}) without receiving any direct
counter-gift.  The ``return'' is not personal reciprocation but the
existence of the system itself---the assurance that if the donor ever
needs blood, others will give in the same spirit.  This is reciprocity
displaced in time and social space: the counter-gift comes from the
\emph{system}, not from any individual partner.

Titmuss's finding that voluntary systems outperform commercial ones
suggests that generalized reciprocity can be more efficient than
balanced exchange when the ``commodity'' (blood) is difficult to verify
for quality.  The formal insight is that when $\depth(a)$ is
observable only imperfectly, the gift logic (which appeals to social
bonds and moral obligation) may elicit higher-depth contributions than
the market logic (which offers monetary compensation for a commodity
whose depth---quality---is hard to price).
\end{example}

% ============================================================
\section{The Void as Non-Contributor}
\label{sec:ch16-void}

Void is the identity element of composition, and in the cooperative context
it represents non-participation: silence before or after a gift.

\begin{theorem}[Void Contributes Nothing---Left (16.1)]
\label{thm:void_joint_left}
For all $b \in \IS$,
\[
  \texttt{jointCompose}(\void, b) = b.
\]
\end{theorem}

\noindent\textit{Proof sketch.}\quad
Direct from the left-identity law of $\compose$: $\compose(\void, b) = b$.
\hfill$\square$

\medskip

\noindent\textbf{Commentary.}
Silence before a gift changes nothing about the gift.  The Maussian
interpretation is that non-participation at the outset imposes no frame on
the exchange; the recipient's contribution stands alone.

\begin{theorem}[Void Contributes Nothing---Right (16.2)]
\label{thm:void_joint_right}
For all $a \in \IS$,
\[
  \texttt{jointCompose}(a, \void) = a.
\]
\end{theorem}

\noindent\textit{Proof sketch.}\quad
From the right-identity law: $\compose(a, \void) = a$.
\hfill$\square$

\medskip

\noindent\textbf{Commentary.}
A gift followed by silence is just the gift.  Reciprocation with void---an
empty counter-gift---leaves the initiator's contribution unaltered.  This
formalizes the asymmetry that Mauss identified: the absence of reciprocation
does not annihilate the gift, but it leaves the social bond incomplete.

\begin{theorem}[Void Symmetry (16.7)]
\label{thm:void_symmetric_left}
For all $a \in \IS$, $\texttt{symmetricJoint}(\void, a)$.
\end{theorem}

\noindent\textit{Proof sketch.}\quad
Both $\compose(\void, a)$ and $\compose(a, \void)$ equal $a$, so
the equation holds trivially.
\hfill$\square$

\medskip

\noindent\textbf{Commentary.}
Void always commutes---silence is the universal symmetric partner.
Non-participation is the one ``gift'' whose order never matters.
This is both mathematically trivial and anthropologically significant:
the absence of contribution creates no ordering conflict.

\begin{theorem}[Void Symmetry---Self (16.8)]
\label{thm:void_symmetric_self}
$\texttt{symmetricJoint}(\void, \void)$.
\end{theorem}

\noindent\textit{Proof sketch.}\quad
$\compose(\void, \void) = \void = \compose(\void, \void)$.
\hfill$\square$

\medskip

\noindent\textbf{Commentary.}
As Wittgenstein remarked, ``Whereof one cannot speak, thereof one must be
silent''---and the order of two silences is immaterial.  This theorem, while
trivial, anchors the degenerate case of the theory.

% ============================================================
\section{Associativity and Multi-Party Cooperation}
\label{sec:ch16-assoc}

The associativity of composition ensures that multi-party cooperation is
well-defined regardless of how sub-coalitions are bracketed.

\begin{theorem}[Associativity of Multi-Party Cooperation (16.3)]
\label{thm:joint_assoc}
For all $a, b, c \in \IS$,
\[
  \texttt{jointCompose}(\texttt{jointCompose}(a, b), c) \;=\;
  \texttt{jointCompose}(a, \texttt{jointCompose}(b, c)).
\]
\end{theorem}

\noindent\textit{Proof sketch.}\quad
Unfolds to $\compose(\compose(a,b), c) = \compose(a, \compose(b,c))$,
which is the associativity axiom of $\compose$.
\hfill$\square$

\medskip

\noindent\textbf{Commentary.}
In a ritual with three participants, it doesn't matter whether $A$ and $B$
cooperate first and then engage $C$, or $A$ engages a pre-formed $B$--$C$
coalition.  The overall composite meaning is the same.  This is a powerful
structural result: it means that multi-party cooperation is intrinsically
well-defined, independent of the contingent order in which sub-coalitions
form.\footnote{In cooperative game theory, associativity of the
characteristic function is often assumed but rarely proved from first
principles.  Here it follows from the algebraic structure of ideatic
space.}

\begin{theorem}[Four-Party Cooperation Associativity (16.18)]
\label{thm:four_party_assoc}
For all $a, b, c, d \in \IS$,
\[
  \texttt{jointCompose}(\texttt{jointCompose}(\texttt{jointCompose}(a, b), c), d) \;=\;
  \texttt{jointCompose}(a, \texttt{jointCompose}(b, \texttt{jointCompose}(c, d))).
\]
\end{theorem}

\noindent\textit{Proof sketch.}\quad
Iterated application of the associativity axiom of $\compose$.
\hfill$\square$

\medskip

\noindent\textbf{Commentary.}
In organizational theory, four-person teams can be sub-divided many ways:
$(AB)(CD)$, $A(B(CD))$, $((AB)C)D$, and so on.  Associativity guarantees
that all bracketings produce the same final product.  The insight extends to
$n$-party cooperation by induction, establishing that the output of a
\texttt{JointProject} is invariant under any re-bracketing of its participant
list.

\begin{theorem}[Empty Project Produces Void (16.9)]
\label{thm:empty_project_void}
\[
  \texttt{output}(\langle [\ ] \rangle) = \void.
\]
\end{theorem}

\noindent\textit{Proof sketch.}\quad
The fold over an empty list returns the initial accumulator $\void$.
\hfill$\square$

\medskip

\noindent\textbf{Commentary.}
An unfunded, unstaffed project yields silence.  This is the organizational
analogue of the void identity: without participants, no meaning is produced.

\begin{theorem}[Singleton Project Preserves Contribution (16.10)]
\label{thm:singleton_project}
For all $a \in \IS$,
\[
  \texttt{output}(\langle [a] \rangle) = a.
\]
\end{theorem}

\noindent\textit{Proof sketch.}\quad
The fold computes $\compose(\void, a) = a$ by left-identity.
\hfill$\square$

\medskip

\noindent\textbf{Commentary.}
A solo artisan's work is their own idea, unmediated by any collaborator.
The project structure does not distort a lone contributor's output.

\begin{example}[Open-Source Software as Cooperative Composition]
\label{ex:open-source}
Open-source software development provides a paradigmatic modern example
of cooperative composition.  Consider the Linux kernel, initiated by
Linus Torvalds in 1991 and developed collaboratively by thousands of
contributors.  Each contribution---a patch, a driver, a subsystem---is
composed with the existing codebase: the project output at time $t$ is
$\texttt{output}(\langle [c_1, c_2, \ldots, c_t] \rangle)$, the
left-fold composition of all contributions in chronological order.

By Theorem~\ref{thm:empty_project_void}, the kernel before any
contribution is $\void$: an empty project produces no output.  By
Theorem~\ref{thm:singleton_project}, Torvalds's initial contribution
$c_1$ is the project in its first form.  Each subsequent contribution
$c_i$ composes with the accumulated result, and by
Theorem~\ref{thm:joint_depth_bound}, the depth of the new output is
bounded by the depth of the existing codebase plus the depth of the new
contribution.

The open-source model is structurally a gift economy: contributors
``give'' their code to the project without receiving direct monetary
compensation.  The counter-gift is the use of the project itself (a
form of generalized reciprocity), the prestige associated with
contribution (a form of symbolic capital in Bourdieu's sense), and the
learning that comes from composing one's ideas with those of expert
collaborators.  Eric Raymond's distinction between the ``cathedral''
(centrally designed software) and the ``bazaar'' (decentralized
open-source development) maps onto our framework as the difference
between a single-composer project and a multi-party joint project with
heterogeneous contributors and emergent compositional structure.
\end{example}

\begin{example}[Academic Co-authorship as Joint Composition]
\label{ex:coauthorship}
Academic co-authorship is a particularly transparent instance of joint
composition.  When two scholars write a paper together, each brings
their individual repertoire---their knowledge, methods, and
theoretical commitments---and the resulting paper is $\compose(a, b)$
where $a$ and $b$ represent the respective contributions.

The order of composition (who is ``first author'') has both social and
structural significance.  In many fields, first authorship signals
primary intellectual contribution---the first author ``sets the frame''
within which the second author's contribution is interpreted.  This
maps precisely onto the asymmetry of $\compose(a, b) \neq \compose(b,
a)$ in the general case (non-symmetric joints, per
Definition~\ref{def:symmetricJoint}).

By Theorem~\ref{thm:joint_depth_bound}, the depth of the co-authored
paper is bounded by $\depth(a) + \depth(b)$: two scholars cannot
produce a paper deeper than the sum of their individual intellectual
capacities.  By Theorem~\ref{thm:trivial_cooperation}, if both
contributions have depth zero, the joint paper has depth zero---
shallow thinking combined with shallow thinking produces only shallow
output.  This formalizes the common academic complaint that
collaborations between mediocre researchers produce mediocre results:
the depth bound is real, and no amount of collaboration can overcome
the limitation of the inputs.
\end{example}

% ============================================================
\section{Resonance Under Cooperation}
\label{sec:ch16-resonance}

Resonance---the fundamental relation of ideatic compatibility---behaves
well under cooperative composition.  This section establishes that
cooperation preserves, propagates, and in some cases strengthens resonance.

\begin{theorem}[Resonant Partners Produce Resonant Outcomes (16.4)]
\label{thm:resonant_partners_resonant_output}
If $\text{resonates}(a_1, a_2)$ and $\text{resonates}(b_1, b_2)$, then
\[
  \text{resonates}\bigl(\texttt{jointCompose}(a_1, b_1),\;
  \texttt{jointCompose}(a_2, b_2)\bigr).
\]
\end{theorem}

\noindent\textit{Proof sketch.}\quad
Follows from the \texttt{resonance\_compose} lemma, which establishes
that resonance is preserved under componentwise composition.
\hfill$\square$

\medskip

\noindent\textbf{Commentary.}
When gift-givers are spiritually aligned---in Mauss's language, when the
\textit{hau} (spirit of the gift) circulates freely---their combined gifts
carry a resonance that mismatched givers cannot achieve.  This theorem
formalizes the intuition that cultural consonance between cooperating parties
produces consonant outcomes.

\begin{theorem}[Reciprocal Exchange Resonance (16.15)]
\label{thm:reciprocal_resonance}
If $\text{resonates}(a, b)$, then
\[
  \text{resonates}\bigl(\texttt{jointCompose}(a, b),\;
  \texttt{jointCompose}(b, a)\bigr).
\]
\end{theorem}

\noindent\textit{Proof sketch.}\quad
Since $\text{resonates}(a, b)$ implies $\text{resonates}(b, a)$ by
symmetry, applying \texttt{resonance\_compose} to the pairs
$(a, b)$ and $(b, a)$ yields the result.
\hfill$\square$

\medskip

\noindent\textbf{Commentary.}
The counter-gift resonates with the gift when giver and receiver are in
spiritual accord.  This is Mauss's central insight formalized: reciprocal
exchange between resonant partners creates a \emph{cycle} of resonant
meaning, regardless of which direction the gift flows.

\begin{theorem}[Self-Resonance of Joint Output (16.11)]
\label{thm:joint_self_resonance}
For all $a, b \in \IS$,
\[
  \text{resonates}\bigl(\texttt{jointCompose}(a, b),\;
  \texttt{jointCompose}(a, b)\bigr).
\]
\end{theorem}

\noindent\textit{Proof sketch.}\quad
Self-resonance is a general property: every idea resonates with itself.
\hfill$\square$

\medskip

\noindent\textbf{Commentary.}
Any completed exchange, however imperfect, achieves internal coherence.
This is the minimal condition for a meaningful cooperative outcome.

\begin{theorem}[Right-Resonance Preservation (16.12)]
\label{thm:joint_resonance_right}
If $\text{resonates}(a, b)$, then for all $c$,
$\text{resonates}\bigl(\texttt{jointCompose}(a, c),\; \texttt{jointCompose}(b, c)\bigr)$.
\end{theorem}

\noindent\textit{Proof sketch.}\quad
Follows from \texttt{resonance\_compose\_right}: resonant agents cooperating
with the same partner produce resonant results.
\hfill$\square$

\medskip

\noindent\textbf{Commentary.}
Durkheim's insight that agents sharing collective consciousness produce
similar ritual outcomes is formalized here: if $a$ and $b$ are culturally
resonant, their joint works with any fixed third party $c$ will resonate.

\begin{theorem}[Left-Resonance Preservation (16.13)]
\label{thm:joint_resonance_left}
If $\text{resonates}(a, b)$, then for all $c$,
$\text{resonates}\bigl(\texttt{jointCompose}(c, a),\; \texttt{jointCompose}(c, b)\bigr)$.
\end{theorem}

\noindent\textit{Proof sketch.}\quad
From \texttt{resonance\_compose\_left}: the same frame applied to resonant
inputs produces resonant outputs.
\hfill$\square$

\medskip

\noindent\textbf{Commentary.}
Bourdieu's concept of habitus finds formal expression here: the same habitus
(frame $c$) applied to resonant inputs produces resonant interpretations.
The frame preserves the resonance of its inputs.

\begin{remark}[The Asymmetry of Composition and Power in Gift Exchange]
\label{rem:asymmetry-power}
The non-commutativity of composition---$\compose(a, b) \neq \compose(b, a)$
in general---has deep implications for the politics of gift exchange.
In Mauss's analysis, the initiator of an exchange exercises a form of
power: by giving first, the initiator establishes the ``frame'' within
which the recipient must respond.  The Kwakiutl potlatch illustrates
this dynamic: the chief who gives the most lavish feast establishes a
standard that rivals must match or exceed, structuring the entire field
of reciprocal obligation.

In compositional terms, the asymmetry means that $\compose(a, b)$ bears
the ``signature'' of $a$ more prominently than that of $b$: the first
argument establishes the context, and the second argument operates
within it.  When the depth of $a$ far exceeds the depth of $b$
($\depth(a) \gg \depth(b)$), the joint output is dominated by $a$'s
contribution---a formal model of the asymmetric gift that creates
subordination rather than solidarity.  Pierre Bourdieu recognized this
dynamic in \textit{The Logic of Practice} (1980): ``the gift that
cannot be reciprocated creates a lasting asymmetry of obligation, which
is another name for power.''

The formal diagnosis of exploitative gift exchange is thus:
$\depth(a) \gg \depth(b)$ combined with social pressure to accept the
composition $\compose(a, b)$ as binding.  Development aid, with its
attendant conditionalities, often takes this form: the donor's
contribution sets a frame that the recipient cannot escape without
rejecting the gift entirely.
\end{remark}

\begin{remark}[Connection to Chapter~14's Auction Theory]
\label{rem:auction-connection}
The cooperative composition framework of this chapter is the
\emph{complement} of the competitive auction framework developed in
Chapter~\ref{ch:auction}.  In an auction, agents compete to have their
ideas selected; in cooperative composition, agents collaborate to
produce a joint idea.  The two frameworks meet in the concept of the
\emph{coalition}: a group of agents who cooperate internally while
competing externally.

Formally, the joint project $\texttt{output}(\langle [a_1, \ldots,
a_k] \rangle)$ produces a composite idea that the coalition then
``enters'' into the marketplace competition of
Chapter~\ref{ch:marketplace} or ``bids'' in the auction of
Chapter~\ref{ch:auction}.  The depth bound of
Theorem~\ref{thm:joint_depth_bound} constrains the coalition's
competitive capacity: no coalition can produce an output deeper than
the sum of its members' depths, which limits the competitive advantage
of collaboration.  This connection between cooperative and competitive
frameworks mirrors the real-world dynamics of firms, political parties,
and cultural movements: internal cooperation generates a collective
product that then competes in the external marketplace of ideas.
\end{remark}

% ============================================================
\section{Depth Bounds on Cooperative Output}
\label{sec:ch16-depth}

Depth measures interpretive complexity.  This section establishes that
cooperative composition respects subadditivity: the complexity of a joint
venture cannot exceed the total intellectual capital invested.

\begin{theorem}[Depth Bound on Joint Output (16.5)]
\label{thm:joint_depth_bound}
For all $a, b \in \IS$,
\[
  \depth\bigl(\texttt{jointCompose}(a, b)\bigr) \;\leq\;
  \depth(a) + \depth(b).
\]
\end{theorem}

\noindent\textit{Proof sketch.}\quad
Direct from the subadditivity of depth under composition:
$\depth(\compose(a,b)) \leq \depth(a) + \depth(b)$.
\hfill$\square$

\medskip

\noindent\textbf{Commentary.}
The complexity of a joint venture cannot exceed the total intellectual
capital invested by both partners.  This is a fundamental resource constraint
on cooperation: two partners of bounded depth can only produce output of
bounded depth.  The analogy to economic production functions is precise---the
output is bounded by the sum of inputs.

\begin{theorem}[Joint With Self Equals ComposeN 2 (16.6)]
\label{thm:joint_self_is_composeN2}
For all $a \in \IS$,
\[
  \texttt{jointCompose}(a, a) = \texttt{composeN}(a, 2).
\]
\end{theorem}

\noindent\textit{Proof sketch.}\quad
By definition, $\texttt{composeN}(a, 2) = \compose(a, \compose(a, \void))
= \compose(a, a)$.
\hfill$\square$

\medskip

\noindent\textbf{Commentary.}
Self-reinforcement through repeated practice is formally identical to
composing an idea with itself.  In ritual theory, the repetition of a
ceremony deepens its effect; here, the depth of the repeated composition is
at most $2 \cdot \depth(a)$ by subadditivity.

\begin{theorem}[Depth Zero Closure Under Cooperation (16.14)]
\label{thm:trivial_cooperation}
If $\depth(a) = 0$ and $\depth(b) = 0$, then
$\depth\bigl(\texttt{jointCompose}(a, b)\bigr) = 0$.
\end{theorem}

\noindent\textit{Proof sketch.}\quad
From the depth-zero closure property of composition.
\hfill$\square$

\medskip

\noindent\textbf{Commentary.}
Shallow discourse combined with shallow discourse can never produce depth.
This is the critical-theory insight that Adorno and Horkheimer might have
appreciated: the culture industry's endless recombination of trivial content
produces only trivial content.\footnote{Theodor Adorno and Max Horkheimer,
\textit{Dialectic of Enlightenment} (1944), argued that mass culture
recombines familiar elements without producing genuine novelty.}

\begin{theorem}[Cooperation Monotonicity for Depth (16.17)]
\label{thm:void_cooperation_depth}
For all $a \in \IS$,
$\depth\bigl(\texttt{jointCompose}(a, \void)\bigr) = \depth(a)$.
\end{theorem}

\noindent\textit{Proof sketch.}\quad
Since $\compose(a, \void) = a$, the depth is preserved.
\hfill$\square$

\medskip

\noindent\textbf{Commentary.}
Free-riders (void contributors) don't add complexity to the joint output,
but they don't reduce it either.  The depth of the project is unchanged
by the presence or absence of non-contributing participants---a result with
clear implications for team management and organizational design.

\begin{remark}[The Spirit of the Gift (\textit{Hau}) as Resonance]
\label{rem:hau-resonance}
Mauss famously invoked the Maori concept of \textit{hau}---the ``spirit
of the gift''---to explain why gifts must be reciprocated.  The
\textit{hau} of an object compels its recipient to make a return; failure
to reciprocate invites spiritual danger.  While Mauss's ethnographic
interpretation has been debated (Marshall Sahlins offered an influential
reinterpretation in \textit{Stone Age Economics}), the structural insight
remains: something in the gift \emph{persists} beyond the physical
transfer, creating an ongoing connection between giver and receiver.

In our framework, the \textit{hau} maps naturally onto \emph{resonance}.
When $a$ gives a gift $g$ to $b$, and $g$ resonates with $b$'s
repertoire ($\text{resonates}(g, r_b)$ for some $r_b \in
\mathbf{R}_b$), the gift creates a compositional connection:
$\compose(g, r_b)$ is a new idea in $b$'s interpretive space that bears
the trace of both the gift and the recipient's prior understanding.
This resonance-trace is the formal analogue of the \textit{hau}: it is
the structural mechanism by which the gift ``stays with'' the giver
even after being transferred.

Theorem~\ref{thm:resonant_partners_resonant_output} formalizes the
positive dynamics of the \textit{hau}: when giver and receiver are in
resonance, their joint compositions produce resonant outcomes.  The
\textit{hau} circulates because resonance is preserved under
composition: the spirit of the gift propagates through the exchange
network as long as the partners remain culturally aligned.  When
resonance breaks down---when the gift fails to resonate with the
recipient's repertoire---the \textit{hau} is lost, and the gift becomes
a mere commodity: an alienable object stripped of its spiritual
connection to the giver.
\end{remark}

% ============================================================
\section{Multi-Party Projects and Iterated Cooperation}
\label{sec:ch16-multi}

We now extend the analysis to multi-party projects, iterated
self-cooperation, and the celebrated gift-exchange cycles of
anthropological theory.

\begin{theorem}[Joint Project Depth Sum Bound (16.16)]
\label{thm:project_depth_bound}
If every participant $a$ in project $p$ satisfies $\depth(a) \leq d$, then
\[
  \texttt{totalDepth}(p) \;\leq\; |p.\text{participants}| \cdot d.
\]
\end{theorem}

\noindent\textit{Proof sketch.}\quad
The total depth is a sum of individual depths, each bounded by $d$.
The sum of $n$ terms each at most $d$ is at most $n \cdot d$.
\hfill$\square$

\medskip

\noindent\textbf{Commentary.}
Total intellectual capital invested constrains what the project can achieve.
This is the project-management analogue of the depth bound: a team of $n$
members, each of bounded capability $d$, produces a project whose total
complexity is at most $n \cdot d$.

\begin{theorem}[Iterated Self-Cooperation Depth Bound (16.19)]
\label{thm:iterated_self_cooperation_depth}
For all $a \in \IS$ and $n \in \mathbb{N}$,
\[
  \depth\bigl(\texttt{composeN}(a, n)\bigr) \;\leq\; n \cdot \depth(a).
\]
\end{theorem}

\noindent\textit{Proof sketch.}\quad
By induction on $n$, using the subadditivity of depth at each step.
\hfill$\square$

\medskip

\noindent\textbf{Commentary.}
Repeating a ritual $n$ times has bounded cumulative depth, growing at most
linearly with repetitions.  This formalizes the intuition that practice
deepens understanding, but at a rate bounded by the practitioner's own
depth.  The logarithmic ``diminishing returns'' often observed in ritual
practice are consistent with this linear upper bound.

\begin{theorem}[Gift Exchange Triangle (16.20)]
\label{thm:gift_triangle_depth}
For all $a, b, c \in \IS$,
\[
  \depth\bigl(\texttt{jointCompose}(\texttt{jointCompose}(a, b), c)\bigr)
  \;\leq\; \depth(a) + \depth(b) + \depth(c).
\]
\end{theorem}

\noindent\textit{Proof sketch.}\quad
Apply the depth bound twice: first to $\compose(a, b)$ and $c$, obtaining
$\depth(\compose(\compose(a,b), c)) \leq \depth(\compose(a,b)) + \depth(c)$,
then to $a$ and $b$, obtaining $\depth(\compose(a,b)) \leq \depth(a) +
\depth(b)$.  Combining by transitivity yields the three-term bound.
\hfill$\square$

\medskip

\noindent\textbf{Commentary.}
This theorem formalizes Malinowski's Kula ring: gifts circulating through a
ring of three partners have bounded total complexity, ensuring the exchange
remains manageable.  The three-term depth bound generalizes to $n$-term
bounds by induction (see Theorem~\ref{thm:iterated_self_cooperation_depth}),
but the triangular case is particularly significant because it corresponds
to the minimal non-trivial cycle of exchange.

% ============================================================
\section{From Mauss to the Digital Commons}
\label{sec:ch16-digital-commons}

The transition from Mauss's analysis of Melanesian shell exchange to
contemporary digital gift economies is not as great a leap as it might
appear.  The structural properties formalized in this chapter---identity,
associativity, resonance preservation, depth subadditivity---hold
regardless of whether the ``gifts'' are shell necklaces, blood, lines
of code, or Wikipedia edits.  This section traces the formal continuities
and examines what changes when gift exchange goes digital.

\subsection{The Creative Commons as Compositional Infrastructure}

The Creative Commons licensing system, founded in 2001, provides a
legal framework for ``keeping-while-giving'' in the digital age.  When
an author releases a work under a Creative Commons license, they give
the work to the public while retaining certain rights (attribution,
non-commercial use, share-alike).  In our framework, this is precisely
Weiner's ``keeping-while-giving'': the idea $a$ enters the public
repertoire via $\operatorname{expand}(\mathbf{R}_{\text{public}}, a)$
without leaving the author's repertoire
$\mathbf{R}_{\text{author}}$.

The ``share-alike'' provision (CC BY-SA) has a particularly elegant
formal interpretation.  It requires that any derivative work---any
composition $\compose(a, b)$ that builds on the original---must itself
be released under the same license.  This is a \emph{compositional
constraint}: the predicate ``is licensed CC BY-SA'' is preserved under
composition.  If $a$ is CC BY-SA and $b$ is composed with it, then
$\compose(a, b)$ must also be CC BY-SA.  The share-alike condition
thus propagates through the compositional network, ensuring that the
gift economy of creative works is self-sustaining: each act of giving
(releasing under CC BY-SA) constrains all future compositions to
remain within the gift economy.

\subsection{Wikipedia as a Multi-Party Joint Project}

Wikipedia, with its hundreds of thousands of active editors, is perhaps
the largest joint project in human history.  Each Wikipedia article is a
joint project in our formal sense: $\texttt{output}(\langle [e_1, e_2,
\ldots, e_n] \rangle)$, where $e_i$ is the $i$-th edit.  By
Theorem~\ref{thm:empty_project_void}, an article before any edits is
$\void$---the empty page.  By Theorem~\ref{thm:singleton_project}, the
first edit $e_1$ establishes the article.  Each subsequent edit composes
with the accumulated content, and the associativity of composition
(Theorem~\ref{thm:joint_assoc}) ensures that the article's meaning is
independent of how we bracket the editorial history.

The ``edit war'' phenomenon---in which opposing editors repeatedly
overwrite each other's contributions---reveals the limits of our
compositional model.  In a strict edit war, the composition oscillates
between $\compose(a, b)$ and $\compose(b, a)$, which are generically
different (non-symmetric joints).  The resolution of edit wars through
consensus mechanisms (talk pages, mediation, administrative fiat)
corresponds to the social process of selecting a canonical composition
order---a choice that our formal framework shows is structurally
consequential, since the output depends on the order of contributions.

\subsection{Cryptocurrency and the Gift-Commodity Boundary}

The emergence of non-fungible tokens (NFTs) and digital collectibles
blurs the boundary between gift and commodity economies in ways that
our framework illuminates.  An artist who ``mints'' an NFT creates a
digital artifact that functions simultaneously as a commodity (it has a
market price) and a gift (it carries the ``aura'' of the creator,
creating a personal connection between artist and collector).  The
\textit{hau} of the NFT is its provenance: the blockchain record of
its creation and subsequent transfers, which is formally analogous to
the compositional history $\compose(a_1, \compose(a_2, \ldots))$ that
accumulates social meaning through circulation.

\paragraph{Extended proof discussion: The gift exchange triangle.}
Theorem~\ref{thm:gift_triangle_depth}, the gift exchange triangle,
merits a more detailed proof analysis.  The theorem states that for
three participants $a, b, c$, the depth of $\compose(\compose(a, b),
c)$ is at most $\depth(a) + \depth(b) + \depth(c)$.  The proof
proceeds by two applications of subadditivity:
\begin{enumerate}
  \item First, $\depth(\compose(\compose(a, b), c)) \le
    \depth(\compose(a, b)) + \depth(c)$ by direct subadditivity.
  \item Second, $\depth(\compose(a, b)) \le \depth(a) + \depth(b)$
    by subadditivity again.
  \item Combining by transitivity of $\le$:
    $\depth(\compose(\compose(a, b), c)) \le
    (\depth(a) + \depth(b)) + \depth(c) =
    \depth(a) + \depth(b) + \depth(c)$.
\end{enumerate}
This telescoping argument generalizes to $n$ participants by induction:
for any joint project $\langle [a_1, \ldots, a_n] \rangle$, the depth
of the output is at most $\sum_{i=1}^n \depth(a_i)$.  The inductive
step uses exactly the same two-step argument, applying subadditivity to
peel off the last participant and invoking the inductive hypothesis for
the remaining $n-1$.

The bound is tight when composition is ``depth-preserving''---i.e., when
$\depth(\compose(a, b)) = \depth(a) + \depth(b)$.  Whether any
non-trivial ideatic space achieves this equality is an open question.
The bound is loose when composition is ``depth-absorbing''---i.e., when
$\depth(\compose(a, b)) \ll \depth(a) + \depth(b)$, as when combining
two redundant ideas produces nothing new.  The gap between the bound and
the actual depth measures the degree of ``compositional redundancy''
between the contributors---a quantity related to mutual information in
information-theoretic terms.

\paragraph{Extended proof discussion: Associativity and its
computational significance.}
The associativity theorem (Theorem~\ref{thm:joint_assoc}) is
algebraically elementary---it follows directly from the associativity
axiom of composition---but its computational and organizational
significance is profound.  In a distributed computing context,
associativity means that a multi-party computation can be parallelized:
instead of computing $\compose(a_1, \compose(a_2, \compose(a_3,
a_4)))$ sequentially (depth $O(n)$), we can compute
$\compose(\compose(a_1, a_2), \compose(a_3, a_4))$ in parallel
(depth $O(\log n)$).  The outputs are identical by associativity.

In organizational terms, this means that hierarchical and flat team
structures produce the same joint output.  A hierarchical organization
computes $\compose(\compose(a_1, a_2), \compose(a_3, a_4))$---pairs
cooperate and then their outputs are composed at the next level.  A
flat organization computes $\compose(a_1, \compose(a_2, \compose(a_3,
a_4)))$---each member contributes sequentially to an accumulating
whole.  Associativity guarantees identical results, which means that
the choice between hierarchical and flat structures is driven by
considerations of communication cost and coordination overhead, not
by the quality of the output.

% ============================================================
\section{Conclusion}
\label{sec:ch16-concl}

This chapter has established the formal foundations of cooperative
composition in ideatic space.  Twenty theorems characterize the algebraic
and metric properties of joint composition, multi-party projects, and
iterated cooperation.

The central results are:
\begin{itemize}
  \item \textbf{Identity:} Void contributes nothing
    (Theorems~\ref{thm:void_joint_left}--\ref{thm:void_joint_right}).
  \item \textbf{Associativity:} Multi-party cooperation is bracket-independent
    (Theorems~\ref{thm:joint_assoc}--\ref{thm:four_party_assoc}).
  \item \textbf{Resonance preservation:} Resonant partners produce resonant
    outcomes (Theorems~\ref{thm:resonant_partners_resonant_output}--\ref{thm:joint_resonance_left}).
  \item \textbf{Depth subadditivity:} Cooperative complexity cannot exceed
    the sum of individual complexities (Theorems~\ref{thm:joint_depth_bound}--\ref{thm:gift_triangle_depth}).
\end{itemize}

Mauss's \textit{Essai sur le don} revealed that gift exchange is a total
social fact.  Our formalization shows that it is also a total \emph{algebraic}
fact: the structure of ideatic space---its identity, associativity, and
subadditivity---completely determines the bounds and resonance properties
of cooperative meaning-making.  In the next chapter, we turn to what happens
when cooperation breaks down: the problem of bargaining over composition order.
